\chapter{第二代总览：两个指纹，两个路标}
\label{ch:gen2-big}

\begin{gentwobox}[第二代 G-TCDC v6 2F-Gear 一句话]
在第一代的基础上，\textbf{多维护一个指纹} $h_{\mathrm{norm}}$；  
切点要同时满足：\textbf{左指纹的路标} 和 \textbf{右指纹的路标}（各占不同 bit 位）。
\end{gentwobox}

\section{为什么要第二代？}

\begin{itemize}[nosep]
  \item G-TCDC 是 \textbf{opt-in} 的另一套 dedup 域（和 FastCDC 仓库不混用）；
  \item 需要在「切点质量」和「扫描速度」之间找新的平衡点；
  \item v6 2F-Gear 的目标：\textbf{切法独立}，但 \textbf{速度接近第一代 Stream}。
\end{itemize}

\section{和第一代的视觉对比}

\begin{figure}[htbp]
\centering
\begin{tikzpicture}[node distance=0.6cm]
  \node[gen1, minimum width=4.2cm, minimum height=1.4cm] (g1) {第一代\\一个 $h$\\一个 mask};
  \node[gen2, minimum width=4.2cm, minimum height=1.4cm, right=1.5cm of g1] (g2)
    {第二代\\$h_{\mathrm{gear}}$ + $h_{\mathrm{norm}}$\\$mask_{lo}$ + $mask_{hi}$};
  \draw[arr] (g1) -- node[above,font=\scriptsize]{进化} (g2);
\end{tikzpicture}
\caption{循环结构不变：仍是每字节 probe 一次；每次多更新一路 hash。}
\end{figure}

\section{第二代切点条件（先背下来，下章揉碎）}

\begin{codebox}[第二代切点（最终版）]
\begin{lstlisting}
左OK  = (h_gear & mask_lo) == 0
右OK  = ((h_norm >> norm_shift) & mask_hi) == 0
if (左OK && 右OK)  在这里切;
\end{lstlisting}
\end{codebox}

\begin{beginnerbox}[三个新名词]
\begin{itemize}[nosep]
  \item \keytermii{mask\_lo}：看 $h_{\mathrm{gear}}$ 的\textbf{低位}；
  \item \keytermii{mask\_hi}：看 $h_{\mathrm{norm}}$ \textbf{移过几位之后}的位；
  \item \keytermii{norm\_shift}：右指纹要右移多少，才能和左指纹「不抢同一位」。
\end{itemize}
\end{beginnerbox}

\section{第二代多一张表}

\begin{figure}[htbp]
\centering
\begin{tikzpicture}[node distance=0.5cm]
  \node[byte] (b) {字节 $b$};
  \node[gen1, below left=0.8cm and 0.2cm of b] (t1) {beta\_table\\（同第一代 gear）};
  \node[gen2, below right=0.8cm and 0.2cm of b] (t2) {norm\_table\\（仓库 seed 派生）};
  \node[gen1, below=0.8cm of t1] (hg) {$h_{\mathrm{gear}}$};
  \node[gen2, below=0.8cm of t2] (hn) {$h_{\mathrm{norm}}$};
  \draw[arr] (b) -- (t1);
  \draw[arr] (b) -- (t2);
  \draw[arr] (t1) -- (hg);
  \draw[arr] (t2) -- (hn);
\end{tikzpicture}
\caption{同一个字节，查两张表，得到两个滚动指纹。}
\end{figure}

滚动公式与第一代\textbf{完全相同}，只是表不同：

\begin{lstlisting}
h_gear = (h_gear << 1) + beta[b]  - beta[b-w]
h_norm = (h_norm << 1) + norm[b]  - norm[b-w]
\end{lstlisting}

\section{模拟示例：同一字节查两表}

\begin{simbox}[$i=3$ 滚动后（未切），为 $i=4$ 准备]
读 \texttt{0x40}、丢 \texttt{0x20}：\\
$h_g = (20 \ll 1) + 66 - 2 = 104$；\quad
$h_n = (8 \ll 1) + 3 - 5 = 16$。\\
下一 probe（$i=4$）两路同时 OK → 切。
\end{simbox}
